% =============================================================================
% 5.8 CALIDAD DEL CÓDIGO Y HERRAMIENTAS DEL FRONTEND
% =============================================================================
\section{Calidad del Código y Herramientas del Frontend}
\label{sec:calidad_fe}

El diseño del frontend incorpora \textbf{garantías de calidad en tiempo de construcción}
que detectan errores antes del despliegue. Esta sección documenta el \textit{tooling} que
las hace cumplir; la estrategia integral de verificación y validación se aborda en el
capítulo~8.

\subsection{Seguridad de tipos con TypeScript estricto}
\label{subsec:typescript}

Todo el código está escrito en \textbf{TypeScript con modo estricto}. El \textit{script}
de construcción ejecuta primero \comando{tsc} (verificación de tipos) y solo si pasa
procede con \comando{vite build}; un error de tipos \textbf{aborta el empaquetado}. Los
tipos de respuesta de los servicios derivan del contrato \texttt{ApiResponse<T>} del
backend (capítulo~\ref{ch:backend}, listado~\ref{lst:apiresponse},
pág.~\pageref{lst:apiresponse}), de modo que un cambio incompatible en la forma de los
datos del servidor se detecta como error de compilación en el cliente.

\subsection{Linting y convenciones (ESLint)}
\label{subsec:eslint}

El proyecto aplica \textbf{ESLint 10} con los \textit{plugins} de React Hooks y React
Refresh, que verifican las reglas de los \textit{hooks} (dependencias completas, llamadas
en el nivel superior) y la compatibilidad con la recarga en caliente. El \textit{script}
\comando{npm run lint} reporta las directivas no usadas y limita las advertencias.

\subsection{Pruebas con Vitest y Testing Library}
\label{subsec:vitest}

Las pruebas unitarias y de integración usan \textbf{Vitest 4} sobre un entorno
\texttt{jsdom}, con \textbf{Testing Library} para interactuar con los componentes como lo
haría un usuario (consultas por rol accesible, eventos de usuario). La tabla~\ref{tab:tooling}
resume la cadena de herramientas de calidad del frontend.

\vspace{0.2cm}
\noindent
\renewcommand{\arraystretch}{1.3}
\fontsize{8.5}{10.2}\selectfont\sffamily
\begin{tabularx}{\textwidth}{|L{3.4cm}|L{3.0cm}|Y|}
\hline
\rowcolor{headerTabla}
\sffamily\textbf{Herramienta} & \sffamily\textbf{Tipo} & \sffamily\textbf{Garantía que aporta} \\
\hline
TypeScript (\texttt{tsc}) & Verificación de tipos & Errores de contrato detectados en compilación. \\ \hline
\rowcolor{grisTecnico}
ESLint & Análisis estático & Reglas de \textit{hooks} y convenciones de código. \\ \hline
Vitest + Testing Library & Pruebas & Comportamiento de componentes y flujos. \\ \hline
\rowcolor{grisTecnico}
Vite (\texttt{build}) & Empaquetado & Optimización y verificación previa al despliegue. \\ \hline
\end{tabularx}
\label{tab:tooling}

\begin{NotaTecnica}
La secuencia \texttt{tsc} $\rightarrow$ \texttt{vite build} convierte la verificación de
tipos en una \textbf{barrera de despliegue}: el empaquetado monolítico del
capítulo~\ref{ch:backend} (sección~\ref{subsec:empaquetado},
pág.~\pageref{subsec:empaquetado}) solo incrusta el \textit{bundle} si el frontend compila
sin errores de tipo. Con esto se cierra el diseño detallado del frontend; el
capítulo~\ref{ch:basedatos} (pág.~\pageref{ch:basedatos}) documenta el motor de datos que
sustenta toda la plataforma.
\end{NotaTecnica}